\section{Verification}
\label{sec:gates}

\subsection{Philosophy}

A check that has never been observed to fail is not evidence. Each gate below
was corrupted deliberately and watched. The two energy checks are reported
together because neither is sufficient alone, and the second exists precisely
because the first proved silent.

\subsection{The adiabatic limit, and an inverted guard}

Driving $k_r \to 0$ must recover hypothesis H6 exactly: zero loss and
$\etast = 1$. On its first execution this test returned $\etast = 0.20$ and a
loss \emph{larger} than the physical case.

The cause was a guard in~\eqref{eq:Rg}. When the argument of the logarithm
falls below unity the line source is outside its validity, and the
implementation returned $\Rg = 0$, intending ``the front has not yet cleared
the hole''. But zero \emph{resistance} is infinite \emph{conductance}: the
expression asserted that the rock was a perfect heat sink, which is the
opposite of what was meant and the most damaging possible failure of that
function. It now returns infinity, and the parameter validation warns rather
than allowing a silent number through. The test passes with zero loss and
$\etast = 1$ to the cyclic convergence tolerance.

\subsection{Why the closure cannot detect a sign error}

Let $q_{\mathrm{loss}}$ be negated in~\eqref{eq:split}. Then
$\qp_{\mathrm{tube}} = q_\Sigma - q_{\mathrm{loss}}$ and the accumulated loss
is $-Q_{\mathrm{loss}}$, so the residual becomes
\begin{equation}
Q_{\mathrm{tube}} - \Delta E - Q_{\mathrm{loss}}^{\,\mathrm{wrong}}
 = \big(\textstyle\sum q_\Sigma - Q_{\mathrm{loss}}\big)
   - \textstyle\sum q_\Sigma - \big(-Q_{\mathrm{loss}}\big) = 0 ,
\end{equation}
identically. The error appears twice with opposite signs and cancels. Measured,
the residual under this corruption is $5\times10^{-15}$ --- machine epsilon.
The closure reads only its own inputs and is therefore blind to the convention
it was written with.

\subsection{A Clausius gate on the loss term}

The independent input the closure lacks is the \emph{state}. Heat must flow
from hot to cold, so $q_{\mathrm{loss}}$ and $(\Tpcm - \Trock)$ must share a
sign at every segment and every step. Accumulating their product,
\begin{equation}
\Sigma \;=\; \sum_{\text{steps}}\sum_{i}
  q_{\mathrm{loss},i}\,\big(T_{p,i} - T_{r,i}\big)\,\Delta z\,\Delta t
  \;=\; \sum K_\ell (\Delta T)^{2} \;\ge\; 0 ,
\end{equation}
which under a reversed convention is $-\sum K_\ell(\Delta T)^2$ and cannot be
positive. The difference from the closure is that $(\Tpcm - \Trock)$ is
evaluated from the marched state rather than from $q_{\mathrm{loss}}$ itself.

\begin{table}[h]
\centering
\begin{tabular}{lrr}
\toprule
corruption & closure & $\Sigma$ \\
\midrule
(none)                           & $2.9\times10^{-16}$ & $+1.64\times10^{11}$ \\
loss omitted from the tube duty  & $5.6\times10^{-1}$ \textbf{fires} & $+1.75\times10^{11}$ silent \\
loss sign reversed               & $5.0\times10^{-15}$ silent & $-3.87\times10^{11}$ \textbf{fires} \\
\bottomrule
\end{tabular}
\caption{Each gate detects exactly what the other cannot. Neither alone would
have been sufficient.}
\end{table}

\subsection{Parameter validation}

Beyond the two energy gates, the case validator checks: the source temperature
against the gradient; the cascade orientation, which must be hot-at-bottom and
whose profile must be in charge-frame indexing; the geometric packing, as a
necessary area condition with a warning when the downcomer must be run
off-centre; the remaining material volume; the fraction of the well whose
melting point lies below the local rock temperature, and so cannot freeze; and
the validity of the line-source form at the requested field age.

